%Para este capítulo se usará la abreviatura "fund".
\chapter{Fundamentos teóricos}
\label{fund}
%Introducción por hacer (pensarlo bien).
\section{Planteamiento del problema}
De ahora en adelante consideraremos fijadas las bases canónicas tanto de $\R^n$ como de $\R$, a no ser que se especifique lo contrario.

Antes de plantear formalmente el problema de la programación lineal, introduzcamos unas definiciones.

\begin{defi}[Vector no negativo]
	Un vector $x=(x_1,\dots,x_n)\in\R^n$ se dice \tbi[vector!no negativo]{no negativo} si todas sus componentes son, como su propio nombre indica, no negativas. Es decir, $x_i\geq 0$ para todo $i\in\{1,\dots,n\}$. Si un vector $x$ es no negativo escribiremos $x\geq 0$.
\end{defi}
Al hilo de esta definición, vamos a generalizar ese concepto que lleva presente en nuestras vidas desde que aquel docente de primaria nos dibujó dos rectas perpendiculares en la pizarra, los cuadrantes.

\begin{defi}[$2^n$--ante positivo]
	Definimos el \tbi{$2^n$--ante positivo} de $\R^n$ como el conjunto de los vectores no negativos de $\R^n$.
	
	Se deja como ejercicio al lector justificar el nombre de $2^n$--ante.
\end{defi}
Dicho esto, ya podemos formular el problema de la programación lineal.
\begin{prob}[Formulación general]
	\label{fund_prob_formulacionGeneral}
	La disciplina de la \tbi[programación!lineal]{programación lineal} estudia procedimientos (implementables en ordenador) para calcular los extremos absolutos de una aplicación lineal $f:\R^n\to\R$ cuando restringimos su dominio una variedad afín de $\R^n$ cortada con el $2^n$--ante positivo.
	
	\begin{figure}[h!]
		\centering
		\label{fund_img_problemaGeneral}
		\includegraphics[scale = 0.75]{img/problemaGeneral}
		\caption{Ilustración del problema de programación lineal.}
	\end{figure}
	En el caso de la ilustración \ref{fund_img_problemaGeneral} suponemos que $f$ es una aplicación lineal definida sobre todo el plano a la cual restringimos el dominio a la variedad afín $\mc{L}$ cortada con el cuadrante positivo (línea azul). Nuestro deber es encontrar los puntos (si los hay) sobre la línea azul en los cuales $f$ alcanza un máximo o un mínimo absoluto.
\end{prob} 
En esta sección nos dedicaremos simplemente a formular de diversas maneras el problema de la programación lineal, observando que todas son equivalentes entre sí. Comencemos pues, sin más dilación nuestro viaje por lo desconocido, un viaje del que probablemente no regresemos.

A continuación vemos una sencilla caracterización de las aplicaciones lineales $f:\R^n\to\R$ a partir de su expresión analítica.
\begin{obs}[Expresión analítica]
	Consideremos una aplicación lineal $f:\R^n\to\R$.
	
	Por ser $f$ una aplicación lineal, tendrá una única matriz asociada $C$. Como salta a la vista, $C$ será fila con $n$ columnas.
	
	Entonces, dado un vector $x\in\R^n$ se tiene que
	\begin{equation}
	\label{fund_eq_lineales}
	f(x)=f(x_1,\dots,x_n)= CX=\begin{pmatrix}
	c_1 & \cdots & c_n
	\end{pmatrix}\begin{pmatrix}
	x_1\\ \cdots \\x_n
	\end{pmatrix}=c_1x_1+\cdots+c_nx_n
	\end{equation}
	De esta forma se concluye que toda aplicación lineal $f:\R^n\to\R$ tiene una expresión analítica de la forma de la ecuación \eqref{fund_eq_lineales}. Recíprocamente, toda función $g:\R^n\to\R$ con una expresión analítica como la de \eqref{fund_eq_lineales} es una ecuación lineal, para demostrarlo, basta con leer como los árabes la ecuación \eqref{fund_eq_lineales}.
\end{obs}
Antes de continuar es importante hacer un inciso acerca de la notación.
\begin{obs}[Notación matricial]
	Siempre que escribamos matrices fila o columna lo haremos con letras minúsculas. Asimismo, siempre que estemos escribiendo una matriz final añadiremos el símbolo de trasposición, para dejarlo claro implícitamente. De esta manera, a las funciones lineales las denotaremos por
	\begin{equation*}
		f(x_1\dots,x_n)=c^tx
	\end{equation*}
	donde $c^t$ es la matriz asociada a $f$ y $x$ es el vector columna con las coordenadas de $x$ respecto de la base canónica.
\end{obs}
\subsection{Forma estándar y formas canónicas}

Sea $f(x_1,\dots,x_n)=c_1x_1+\cdots+c_nx_n=c^tx$ una función lineal  a la que a partir de ahora llamaremos \tbi[función!objetivo]{función objetivo}. A la matriz $c$ se la denomina \tbi[vector!de coeficientes de la función objetivo]{vector de coeficientes de la función objetivo}, mientras que $x$ recibe el nombre de \tbi[vector!de variables de decisión]{vector de variables de decisión}.

Consideremos asimismo la variedad afín de $\R^n$ dada por el conjunto de soluciones al sistema no necesariamente homogéneo $Ax=b$. Escrito de forma desarrollada
\begin{equation*}
\begin{array}{cc}
\left\{\begin{array}{cccc}
a_{11}x_1&+\cdots+&a_{1n}x_n&=b_1\\
\vdots &\ddots &\vdots &\vdots\\
a_{1m}x_1&+\cdots+&a_{mn}x_n&=b_m
\end{array}\right.\qquad&\qquad\begin{pmatrix}
a_{11} & \cdots & a_{1n}\\
\vdots & \ddots & \vdots\\
a_{m1} & \cdots & a_{mn}
\end{pmatrix}\begin{pmatrix}
x_1\\
\vdots\\
x_n
\end{pmatrix}=\begin{pmatrix}
b_1\\
\vdots\\
b_m
\end{pmatrix}
\end{array}
\end{equation*}
A la matriz $A\in\mathfrak{M}_{m\times n}(\R)$ usualmente la llamaremos \tbi[matriz!de coeficientes de las restricciones]{matriz de coeficientes de las restricciones}. Asimismo el vector $b$ será conocido como \tbi[vector!de términos independientes]{vector de términos independientes}.
\begin{prob}[Forma estándar]
El problema de programación lineal en \tbi[forma!estándar]{forma estándar} consiste en hallar los vectores $x\geq 0$ tales que sean solución de $Ax=b$ y además sean mínimos absolutos de la función $f$. Usualmente se escribe de forma sintética (aunque no lo usaremos mucho)
\begin{equation*}
\begin{array}{c}
\min c^tx\\
\text{Sujeto a:}\qquad Ax=b,\qquad x\geq 0
\end{array}
\end{equation*}
Nótese que la formulación del problema en forma estándar es exactamente la misma que la formulación general \ref{fund_prob_formulacionGeneral} con la única diferencia de que ya no buscamos extremos absolutos en general, sino únicamente mínimos. 
\end{prob}
\section{Programación lineal y conjuntos convexos}
\subsection{Conjuntos convexos. Definición y terminología}
\subsection{Puntos extremos}
\subsection{Direcciones extremas}
\section{Teoremas fundamentales}
\section{Algoritmo del Símplex}
\subsection{Fundamentos teóricos}
\subsection{Recapitulación}
\section{Ejercicios resueltos}
\begin{exerc}
	contenidos...
\end{exerc}
\begin{solu}
	contenidos...
\end{solu}